% code_festival_2017_qualc B - Similar Arrays の手書きメモの再現（振り返りの文章は thinking.md）
% ビルド: tools/build_thinking.sh other/code-festival-2017-qualc/b --preview
\documentclass[a4paper]{article}
\usepackage[margin=1.2cm]{geometry}
\usepackage{handmemo}
\pagestyle{empty}

\begin{document}

% ---------------------------------------------------------------- 1 枚目
\begin{memopage}{20260925\_174806.jpg}
  \hw(1.6,1.7){B --- Similar Arrays}

  \hw(1.0,3.4){$\forall i \in [1,N]),\ |x_i - y_i| \leq 1$}
  \hw(9.0,2.6){$x, y$}
  \hw(8.8,3.4){似ている}

  \hw(2.1,4.9){$-1 \leq x_i - y_i \leq 1$}
  \hw(3.1,5.9){$x_i - y_i =$}
  \hwm(5.2,5.9){\left\{\begin{array}{l} 0 \\ \pm 1 \end{array}\right.}

  \hw(0.7,7.7){b で積 : odd $\leftarrow$}
  \hw(4.8,7.7){$\exists\,2$}
  \hwunread(5.9,7.7){1 字（「数」か）}
  \hwline(3.1,8.3)(3.35,8.55) \hwline(3.35,8.55)(3.9,8.15)
  \hw(5.3,8.5){$\uparrow$}
  \hw(4.3,9.4){$\neg\exists$ : 全 odd}
  \redpen(9.6,7.2){積が偶数 ⇔ どれか 1 項が偶数（$\exists$）。\\直接は数えにくいので、余事象\\「全部奇数」（$\neg\exists$）を数える}
  \redarrow(9.5,8.4)(8.1,9.3)

  \hw(5.1,11.3){全:}
  \hw(3.9,12.3){各項で 3}
  \hwunread(6.25,12.3){2 字}
  \hw(8.35,12.3){\small ずつ}
  \hwline(9.0,10.5)(9.1,14.0) \hwline(9.1,14.0)(11.9,13.95)
  \hw(9.4,10.9){全 odd}
  \hw(9.3,12.2){偶の個数 m}
  \hwcircle(10.7,12.2)(0.32,0.38)
  \hwline(9.4,12.65)(12.0,12.6)
  \hwm(10.3,13.25){2^{m}}
  \hwm(5.9,13.3){3^{n}}
  \hwcircle(6.35,13.05)(0.28,0.4)
  \redpen(12.4,11.3){偶数の $A_i$ は\\$A_i \pm 1$ の 2 通り、\\奇数の $A_i$ は\\そのままの 1 通り}

  \hw(1.2,11.8){2}
  \hw(1.2,13.0){2\ \ 3}
  \hw(2.8,13.3){\normalsize m=1}
  \hw(1.45,13.85){\normalsize $\downarrow$}
  \hw(1.3,14.3){1\ 3}
  \hwline(1.2,14.75)(2.2,14.75)

  \hwm(5.4,15.2){3^{n} - 2^{m}}
  \redbox(5.1,14.4)(8.6,16.0)
  \redpen(8.9,15.2){これが答え（1 回目の提出で AC）}

  \hwm(6.4,16.7){3^{2} - 2^{1} = 9 - 2 = 7}
  \hw(10.3,16.0){\normalsize (}

  \hw(1.3,17.9){3}
  \hw(1.2,18.9){3\ 3\ 3}
  \hwm(5.0,18.5){3^{3} -}
  \hwline(6.0,18.3)(6.35,18.7) \hwline(6.0,18.7)(6.35,18.3)
  \hwm(6.45,18.5){2^{0} = 26}
  \redpen(1.2,20.6){$m=0$（全部奇数）の端も検算していた\\→ $3^3 - 2^0 = 26$ はサンプルと一致}
\end{memopage}

% ---------------------------------------------------------------- 図
\clearpage
\section*{図: $A = (2, 3)$ で数える}
全体は各項 3 通り、全部奇数（赤丸）は偶数の項だけ 2 通り・奇数の項は 1 通り。

\begin{center}
\begin{tikzpicture}[x=1.3cm,y=1cm]
  % 例: A = (2, 3)
  \node[anchor=east] at (-0.3,1) {$A_1 = 2$};
  \node[anchor=east] at (-0.3,0) {$A_2 = 3$};
  \foreach \v/\k in {1/0,2/1,3/2} {
    \draw (\k,1) circle (0.28) node {\small $\v$};
  }
  \foreach \v/\k in {2/0,3/1,4/2} {
    \draw (\k,0) circle (0.28) node {\small $\v$};
  }
  \draw[redpen,line width=1.1pt] (0,1) circle (0.36);
  \draw[redpen,line width=1.1pt] (2,1) circle (0.36);
  \draw[redpen,line width=1.1pt] (1,0) circle (0.36);
  \node[anchor=west] at (2.7,1) {\small 3 通り（奇数は \textcolor{redpen}{2 通り}）};
  \node[anchor=west] at (2.7,0) {\small 3 通り（奇数は \textcolor{redpen}{1 通り}）};
  \node[anchor=west] at (2.7,-1) {\small 全体 $3 \times 3 = 9$、全部奇数 \textcolor{redpen}{$2 \times 1 = 2$} → 答え $9 - 2 = 7$};
\end{tikzpicture}
\end{center}

\end{document}
